% !Mode:: "TeX:UTF-8"

\chapter[轨迹预测]{考虑协方差特性的半解析预测方法}[Trajectory Prediction]
半解析预测方法结合了数值法精度高和解析法速度快的优势，对交会任务的大规模初值搜索提供巨大便利。尽管如此，在多种摄动及空间不确定因素下，半解析解的特性对不确定性的统计考虑对于以其为基础的轨道确定任务至关重要。本章将充分吸收各类半解析解结构优势，进一步考虑统计特性，开发更适合实时轨道确定的半解析解。

\section{轨道特性}[Orbit Characteristics]

\subsection{不同轨道类型}
\subsubsection{拱线进动}
\subsubsection{相对距离变化}
\subsection{非线性姿态动力学}
\subsection{大偏心率的几个视角}
\subsection{轨道倾角的影响}


\section{应用案例}

\subsection{同步转移轨道}
\subsection{闪电轨道}
$
3.4.1 闪电轨道共振情形J_2,2
3.4.2 苔原轨道共振情形
3.4.3 考虑大气阻力特性
3.4.4 用于载人登月的过渡轨道
3.4.5 用于发射DRO的过渡轨道
PROBA-3大椭圆高精度编队
$
\subsection{实际数据验证}

\section{本章小结}
本章结合力学模型分析大椭圆轨道动力学特性，总结出轨道递推的核心要点，计算速度和精度。提出一种半解析方法，用于轨迹预测，该方法计算量与DSST相比，与HEOSAT相比计算速度更加，对于。对于协方差与状态转移矩阵，状态转移张量的考虑重要性，进一步给出了考虑状态转移张量的方法。
针对非线性相对运动特性，给出一种解析表达式，该表达式形成一个时变流形。这是后文最优交会规划的条件之一。


% Local Variables:
% TeX-master: "../thesis"
% TeX-engine: xetex
% End:
